\begin{center}{\LARGE\bfseries 参考答案与解析}\par\small 原创模拟试卷 · 评分建议非官方标准\end{center}
\SectionTitle{客观题答案速查}
\begin{tabularx}{\linewidth}{@{}l*{8}{>{\centering\arraybackslash}X}@{}}
\hline
题号 & 1 & 2 & 3 & 4 & 5 & 6 & 7 & 8 \\
答案 & B & C & B & C & D & B & C & A \\ \hline
题号 & 9 & 10 & 11 & 12 & 13 & 14 & & \\
答案 & ABD & ABC & ABC & 28 & 4 & 72 & & \\ \hline
\end{tabularx}
\SectionTitle{一、单项选择题答案}
\Question{\textbf{B}。两个集合的公共部分为 $[0,3]$。}
\Question{\textbf{C}。$\dfrac{1+2i}{1-i}=\dfrac{(1+2i)(1+i)}2=-\dfrac12+\dfrac32i$，虚部是实数 $\dfrac32$。}
\Question{\textbf{B}。通项为 $\binom6k x^{6-2k}$，令 $6-2k=2$，得 $k=2$，系数为15。}
\Question{\textbf{C}。$\cos\alpha=\dfrac45$，故 $\sin2\alpha=2\cdot\dfrac35\cdot\dfrac45=\dfrac{24}{25}$。}
\Question{\textbf{D}。$\boldsymbol a\cdot\boldsymbol b=t-2=0$，故 $t=2$。}
\Question{\textbf{B}。均值为8，方差为 $(16+4+0+4+16)/5=8$。}
\Question{\textbf{C}。正态分布关于2对称，故 $P(X>5)=P(X<-1)=p$，所求为 $\dfrac12-p$。}
\Question{\textbf{A}。必有 $a\geq0$。此时方程等价于 $(x-1)^2=1+a$ 或 $(x-1)^2=1-a$。要有四个不同实根，两个右端须均为正且不同，恰为 $0<a<1$。}
\SectionTitle{二、多项选择题答案}
\Question{\textbf{ABD}。周期为 $2\pi/2=\pi$；$x=\pi/3$ 处相位为 $\pi/2$，是对称轴；在所给区间上相位从 $\pi/2$ 增至 $3\pi/2$，函数单调递减；代入 $x=\pi/12$ 得0。}
\Question{\textbf{ABC}。$X\sim B(3,1/2)$，故 $E(X)=3/2$，$D(X)=3/4$。$P(X\geq2)=(3+1)/8=1/2$，$P(X=1)=3/8$。}
\Question{\textbf{ABC}。由基本不等式，$x+1/x\geq2$，当且仅当 $x=1$ 取等号。$f'(x)=1-1/x^2<0$ 在 $(0,1)$ 成立。$f(x)=3$ 等价于 $x^2-3x+1=0$，两根均为正且乘积为1。又 $1/(x+y)<1/x+1/y$，所以D中的不等号方向错误。}
\SectionTitle{三、填空题答案}
\Question{\textbf{28}。公差 $d=(19-7)/4=3$，故 $a_{10}=19+3d=28$。}
\Question{\textbf{4}。圆心到直线的距离为 $c/\sqrt2$，等于半径 $2\sqrt2$，故 $c=4$。}
\Question{\textbf{72}。总排列数为 $5!$，甲乙相邻的排列数为 $2\cdot4!$，相减得72。}
\newpage
\SectionTitle{四、解答题解析与建议评分}
\Question{解答与评分建议}
\begin{Subquestions}
\item \Points{4}正方形中 $BC\perp AB$，且 $SA\perp$ 平面 $ABCD$，故 $BC\perp SA$。由于 $AB,SA$ 是平面 $SAB$ 内两条相交直线，故 $BC\perp$ 平面 $SAB$。
\item \Points{4}以 $A$ 为原点，$AB,AD,AS$ 分别为 $x,y,z$ 轴建立直角坐标系，则各点坐标为：
\[ A(0,0,0),\quad B(2,0,0),\quad C(2,2,0),\quad D(0,2,0),\quad S(0,0,2). \]
平面 $SBC$ 的方程为 $x+z=2$，所以距离为 $2/\sqrt2=\sqrt2$。
\item \Points{5}该平面的法向量可取 $\boldsymbol n=(1,0,1)$，$\overrightarrow{SD}=(0,2,-2)$。设线面角为 $\theta\in[0,\pi/2]$，则
\[
\sin\theta=\dfrac{|\overrightarrow{SD}\cdot\boldsymbol n|}{|\overrightarrow{SD}|\,|\boldsymbol n|}=\dfrac2{2\sqrt2\cdot\sqrt2}=\dfrac12.
\]
\end{Subquestions}
\Question{解答与评分建议}
\begin{Subquestions}
\item \Points{4}两台均不正确的概率为 $(1/5)(1/4)=1/20$，故至少一台正确的概率为 $19/20$。
\item \Points{4}两台均正确的概率为 $(4/5)(3/4)=3/5$，所求条件概率为 $(3/5)/(19/20)=12/19$。
\item \Points{7}设3个样品中甲检测错误的个数为 $Y$，则 $Y\sim B(3,1/5)$，$T=3+Y$。分布列为 $P(T=3)=64/125$，$P(T=4)=48/125$，$P(T=5)=12/125$，$P(T=6)=1/125$。因此 $E(T)=3+3/5=18/5$。乙的正确率不影响检测次数，因为每个样品至多由乙补检一次。
\end{Subquestions}
\Question{解答与评分建议}
\begin{Subquestions}
\item \Points{3}由焦点得 $c=1$，由点 $(0,\sqrt3)$ 得 $b^2=3$，故 $a^2=b^2+c^2=4$，椭圆为 $x^2/4+y^2/3=1$。
\item \Points{6}设直线单位方向向量为 $(p,q)$，其中 $p^2+q^2=1$，以 $(x,y)=(1+pt,qt)$ 表示直线上点。代入椭圆得 $(3+q^2)t^2+6pt-9=0$。设两根为 $t_-<0<t_+$，则 $t_-t_+=-9/(3+q^2)$，判别式为 $36p^2+36(3+q^2)=144$，故 $t_+-t_-=12/(3+q^2)$。两段焦点距离分别为 $-t_-,t_+$，所以倒数和为 $(t_+-t_-)/(-t_-t_+)=4/3$。
\item \Points{6}弦长 $|AB|=12/(3+q^2)$，原点到直线的距离为 $|q|$，故面积 $K=6|q|/(3+q^2)$。令 $z=|q|\in[0,1]$，则 $3/2-K=3(z-1)(z-3)/(2(3+z^2))\geq0$。当且仅当 $z=1$ 取等号，此时直线为 $x=1$，面积最大值为 $3/2$。用参数式已包含竖直直线，无斜率不存在的遗漏。
\end{Subquestions}
\newpage
\Question{解答与评分建议}
\begin{Subquestions}
\item \Points{4}将递推式除以 $2^n$，得 $a_{n+1}/2^n=a_n/2^{n-1}+1$。因此 $a_n/2^{n-1}$ 是首项为1、公差为1的等差数列，$a_n=n2^{n-1}$。
\item \Points{5}$S_n=1+2\cdot2+3\cdot2^2+\cdots+n2^{n-1}$，错位相减得 $S_n=(n-1)2^n+1$。
\item \Points{8}取 $n=2$ 得 $\lambda\leq S_2/8=5/8$。下面证明 $S_n\geq5n^3/8$。$n=1$ 时成立，$n=2$ 时取等号。先证明对 $n\geq2$，$8(n+1)2^n\geq5(3n^2+3n+1)$：$n=2$ 时左侧96、右侧95；若第 $n$ 项成立，则第 $n+1$ 项的左侧严格大于第 $n$ 项左侧的2倍，而 $2(3n^2+3n+1)-[3(n+1)^2+3(n+1)+1]=3n^2-3n-5>0$，故结论由归纳法成立。于是 $S_{n+1}-S_n=(n+1)2^n\geq(5/8)[(n+1)^3-n^3]$。从 $n=2$ 的等号累加，即得对所有 $n\geq2$ 的目标不等式。故所求为 $\lambda\in(-\infty,5/8]$。
\end{Subquestions}
\Question{解答与评分建议}
\begin{Subquestions}
\item \Points{4}$f_1'(x)=\ln x$，函数在 $(0,1)$ 上递减，在 $(1,+\infty)$ 上递增，故最小值为 $f_1(1)=0$。
\item \Points{6}$f_a'(x)=\ln x+1-a$，唯一驻点为 $x_0=e^{a-1}$。函数先减后增，最小值为 $1-e^{a-1}$；且 $\lim_{x\to0^+}f_a(x)=1$，$\lim_{x\to+\infty}f_a(x)=+\infty$。因此 $a<1$ 时无零点，$a=1$ 时有一个零点1，$a>1$ 时有两个不同零点。
\item \Points{7}由（2）知 $a>1$。又 $f_a(1)=1-a<0$，所以 $0<u<1<v$。令 $h(x)=\ln x+1/x$，由零点方程得 $h(u)=h(v)=a$。$h'(x)=(x-1)/x^2$，故 $h$ 在 $(1,+\infty)$ 上严格递增。对 $t>1$，令 $g(t)=t-1/t-2\ln t$，有 $g(1)=0$，$g'(t)=(t-1)^2/t^2>0$，故 $g(t)>0$。取 $t=1/u>1$，可得 $h(u)-h(1/u)=t-1/t-2\ln t>0$。于是 $h(v)>h(1/u)$，由递增性得 $v>1/u$，即 $uv>1$。再由基本不等式，$u+v\geq2\sqrt{uv}>2$。
\end{Subquestions}
